
%%%% Rego

%% What will this section describe?
Rego Policy Language is the specification language for Open Policy Agent (OPA).
Rego is similar to Datalog TODO in its structure.
It is a declarative language that focuses heavily on its readability and ensuring unambiguity of queries.
On the other hand, it provides powerful support for working with nested documents.
Like other applications which support declarative query languages, OPA is able to optimize queries to improve performance.
This section describes the key characteristicts of the Rego policy language, based on the official documentation of Open Policy Agent (OPA) TODO.

%%% What are the key features of rego policies
% \subsection{Input and Data}

For Rego, query is the input in the form of structured document, usually in JSON format.
To evaluate it properly, the OPA needs to have access to the internal data, which can store, for example, forbidden servers, allowed actions for users, list of administrators, etc.
...

\subsection{Rules}

A Rego policy consists of three parts: package specification, imports, and rules.
Package specification is a manadory line, which specifies the package for current policy.
Imports are voluntary declarations of other packages, that the current policy is dependent on, i.e., it some of its functions.

%%% What is the input / data
%%% How do the rules work

%%%% END Rego %%%%
